\Scene{4}[The Forest of Arden]

\enter{\textsc{Duke Senior, Amiens, Jaques, Orlando, Oliver,} and \textsc{Celia} as \textsc{Aliena}}


\begin{verse_speech}[Duke Senior]
Dost thou believe, Orlando, that the boy\\
Can do all this that he hath promised?
\end{verse_speech}

\begin{verse_speech}[Orlando]
I sometimes do believe and sometimes do not,\\
As those that fear they hope, and know they fear.
\end{verse_speech}

\enter{\textsc{Rosalind as Ganymede, Silvius,} and \textsc{Phoebe}}


\begin{verse_speech}[Rosalind as Ganymede]  
Patience once more whiles our compact is urged.\\
\stage{To \textsc{Duke}} You say, if I bring in your Rosalind,\\
You will bestow her on Orlando here?
\end{verse_speech}

\verseline[Duke Senior]{That would I, had I kingdoms to give with her.}

\begin{verse_speech}[Rosalind as Ganymede]  
\stage{To \textsc{Orlando}} And you say you will have her when I bring her?
\end{verse_speech}

\verseline[Orlando]{That would I, were I of all kingdoms king.}

\begin{a4}
	\pagebreak[4]
\end{a4}

\begin{verse_speech}[Rosalind as Ganymede]  
\stage{To \textsc{Phoebe}} You say you'll marry me if I be willing?
\end{verse_speech}

\verseline[Phoebe]{That will I, should I die the hour after.}

\begin{verse_speech}[Rosalind as Ganymede]  
But if you do refuse to marry me,\\
You'll give yourself to this most faithful shepherd?
\end{verse_speech}

\verseline[Phoebe]{So is the bargain.}

\begin{verse_speech}[Rosalind as Ganymede]  
\stage{To \textsc{Silvius}} You say that you'll have Phoebe if she will?
\end{verse_speech}

\verseline[Silvius]{Though to have her and death were both one thing.}

\begin{verse_speech}[Rosalind as Ganymede]  
I have promised to make all this matter even.\\
Keep you your word, O duke, to give your daughter,—\\
You yours, Orlando, to receive his daughter.—\\
Keep you your word, Phoebe, that you'll marry me,\\
Or else, refusing me, to wed this shepherd.—\\
Keep your word, Silvius, that you'll marry her\\
If she refuse me. And from hence I go\\
To make these doubts all even.
\end{verse_speech}


\exit{\textsc{Rosalind} and \textsc{Celia}}

\begin{verse_speech}[Duke Senior]
I do remember in this shepherd boy\\
Some lively touches of my daughter's favour.
\end{verse_speech}

\begin{verse_speech}[Orlando]
My lord, the first time that I ever saw him\\
Methought he was a brother to your daughter.\\
But, my good lord, this boy is forest-born\\
And hath been tutored in the rudiments\\
Of many desperate studies by his uncle,\\
Whom he reports to be a great magician\\
Obscured in the circle of this forest.
\end{verse_speech}

\enter{\textsc{Touchstone} and \textsc{Audrey}}


\verseline[Jaques]{There is sure another flood toward, and these couples are coming to the ark. Here comes a pair of very strange beasts, which in all tongues are called fools.}

\verseline[Touchstone]{Salutation and greeting to you all.}

\begin{prose_speech}[Jaques]
\stage{To \textsc{Duke}}  Good my lord, bid him welcome. This is the motley-minded gentleman that I have so often met in the forest. He hath been a courtier, he swears.
\end{prose_speech}

\verseline[Touchstone]{If any man doubt that, let him put me to my purgation. I have trod a measure. I have flattered a lady. I have been politic with my friend, smooth with mine enemy. I have undone three tailors. I have had four quarrels, and like to have fought one.}

\verseline[Jaques]{And how was that ta'en up?}

\verseline[Touchstone]{Faith, we met and found the quarrel was upon the seventh cause.}

\verseline[Jaques]{How <seventh cause>?—Good my lord, like this fellow.}

\verseline[Duke Senior]{I like him very well.}

\begin{prose_speech}[Touchstone]
  God 'ild you, sir. I desire you of the like. I press in here, sir, amongst the rest of the country copulatives, to swear and to forswear, according as marriage binds and blood breaks. A poor virgin, sir, an ill-favoured thing, sir, but mine own. A poor humour of mine, sir, to take that that no man else will. Rich honesty dwells like a miser, sir, in a poor house, as your pearl in your foul oyster.
\end{prose_speech}

  \begin{colorbigpic}
	[\basicscale]
	{5iv_mine}
	[An ill-favoured thing, sir, but mine own]
\end{colorbigpic}

\verseline[Duke Senior]{By my faith, he is very swift and sententious.}

\verseline[Touchstone]{According to the fool's bolt, sir, and such dulcet diseases.}

\verseline[Jaques]{But for the seventh cause. How did you find the quarrel on the seventh cause?}

\begin{prose_speech}[Touchstone]
 Upon a lie seven times removed.—Bear your body more seeming, Audrey.—As thus, sir: I did dislike the cut of a certain courtier's beard. He sent me word if I said his beard was not cut well, he was in the mind it was. This is called <the retort courteous.> If I sent him word again it was not well cut, he would send me word he cut it to please himself. This is called <the quip modest.> If again it was not well cut, he disabled my judgment. This is called <the reply churlish.> If again it was not well cut, he would answer I spake not true. This is called <the reproof valiant.> If again it was not well cut, he would say I lie. This is called <the countercheck quarrelsome,< and so to >the lie circumstantial,> and <the lie direct.>
\end{prose_speech}

\verseline[Jaques]{And how oft did you say his beard was not well cut?}

\verseline[Touchstone]{I durst go no further than the lie circumstantial, nor he durst not give me the lie direct, and so we measured swords and parted.}

\verseline[Jaques]{Can you nominate in order now the degrees of the lie?}

\begin{prose_speech}[Touchstone]
  O sir, we quarrel in print, by the book, as you have books for good manners. I will name you the degrees: the first, <the retort courteous>; the second, <the quip modest>; the third, <the reply churlish<; the fourth, >the reproof valiant>; the fifth, <the countercheck quarrelsome>; the sixth, <the lie with circumstance>; the seventh, <the lie direct.> All these you may avoid but the lie direct, and you may avoid that too with an <if.> I knew when seven justices could not take up a quarrel, but when the parties were met themselves, one of them thought but of an <if,> as: <If you said so, then I said so.> And they shook hands and swore brothers. Your <if> is the only peacemaker: much virtue in <if.>
\end{prose_speech}

\begin{verse_speech}[Jaques]
\stage{To \textsc{Duke}}  Is not this a rare fellow, my lord? He's as good at anything and yet a fool.
\end{verse_speech}

\verseline[Duke Senior]{He uses his folly like a stalking-horse, and under the presentation of that he shoots his wit.}

\enter{\textsc{Hymen, Rosalind,} and \textsc{Celia}}

\stage{Music}

\begin{verse_speech}[Hymen]
		\indentpattern{00100100}
	\begin{verse}
	\begin{patverse}
	Then is there mirth in heaven\\
	When earthly things made even\\
	   Atone together.\\
	Good duke, receive thy daughter.\\
	Hymen from heaven brought her,\\
	   Yea, brought her hither,\\
	That thou mightst join her hand with his,\\
	Whose heart within his bosom is.\\
		\end{patverse}
	\end{verse}
\end{verse_speech}

\begin{verse_speech}[Rosalind]
\stage{To \textsc{Duke}}
To you I give myself, for I am yours.
\stage{To \textsc{Orlando}} To you I give myself, for I am yours.
\end{verse_speech}

\verseline[Duke Senior]{If there be truth in sight, you are my daughter.}

\verseline[Orlando]{If there be truth in sight, you are my Rosalind.}

\begin{verse_speech}[Phoebe]
If sight and shape be true,\\
Why then, my love adieu.
\end{verse_speech}

\begin{verse_speech}[Rosalind]
\stage{To \textsc{Duke}} I'll have no father, if you be not he.\\
\stage{To \textsc{Orlando}} I'll have no husband, if you be not he,\\
\stage{To \textsc{Phoebe}} Nor ne'er wed woman, if you be not she.
\end{verse_speech}

\begin{verse_speech}[Hymen]
	\indentpattern{001001}
	\begin{verse}
	\begin{patverse}
	Peace, ho! I bar confusion.\\
	'Tis I must make conclusion\\
	   Of these most strange events.\\
	Here's eight that must take hands\\
	To join in Hymen's bands,\\
	   If truth holds true contents.\\
	\end{patverse}
	\end{verse}
\stage{To \textsc{Rosalind} and \textsc{Orlando}}
	You and you no cross shall part.
\stage{To \textsc{Celia} and \textsc{Oliver}}
	You and you are heart in heart.
\stage{To \textsc{Phoebe}}
	You to his love must accord\\
	Or have a woman to your lord.
\stage{To \textsc{Audrey} and \textsc{Touchstone}}
	You and you are sure together\\
	As the winter to foul weather.
\stage{To \textsc{All}}
	Whiles a wedlock hymn we sing,\\
	Feed yourselves with questioning,\\
	That reason wonder may diminish\\
	How thus we met, and these things finish.

\stage{Song}
\indentpattern{010100}

	\begin{verse}
	\begin{patverse}
	Wedding is great Juno's crown,\\
	   O blessèd bond of board and bed.\\
	'Tis Hymen peoples every town.\\
	   High wedlock then be honoured.\\
	Honour, high honour, and renown\\
	To Hymen, god of every town.
	\end{patverse}
	\end{verse}
\end{verse_speech}

\begin{verse_speech}[Duke Senior]
\stage{To \textsc{Celia}}
O my dear niece, welcome thou art to me,\\
Even daughter, welcome in no less degree.
\end{verse_speech}

\begin{verse_speech}[Phoebe]
\stage{To \textsc{Silvius}}
I will not eat my word. Now thou art mine,\\
Thy faith my fancy to thee doth combine.
\end{verse_speech}

\enter{\textsc{Second Brother, Jaques de Boys}}


\begin{verse_speech}[Second Brother]
Let me have audience for a word or two.\\
I am the second son of old Sir Rowland,\\
That bring these tidings to this fair assembly.\\
Duke Frederick, hearing how that every day\\
Men of great worth resorted to this forest,\\
Addressed a mighty power, which were on foot\\
In his own conduct, purposely to take\\
His brother here and put him to the sword;\\
And to the skirts of this wild wood he came,\\
Where, meeting with an old religious man,\\
After some question with him, was converted\\
Both from his enterprise and from the world,\\
His crown bequeathing to his banished brother,\\
And all their lands restored to them again\\
That were with him exiled. This to be true\\
I do engage my life.
\end{verse_speech}

\begin{verse_speech}[Duke Senior]
\hspace{\widthof{I do engage my life.}}Welcome, young man.\\
Thou offer'st fairly to thy brothers' wedding:\\
To one his lands withheld, and to the other\\
A land itself at large, a potent dukedom.—\\
First, in this forest let us do those ends\\
That here were well begun and well begot,\\
And, after, every of this happy number\\
That have endured shrewd days and nights with us\\
Shall share the good of our returned fortune\\
According to the measure of their states.\\
Meantime, forget this new-fall'n dignity,\\
And fall into our rustic revelry.—\\
Play, music.—And you brides and bridegrooms all,\\
With measure heaped in joy to th' measures fall.
\end{verse_speech}

\begin{verse_speech}[Jaques]
\stage{To \textsc{Second Brother}}
Sir, by your patience: if I heard you rightly,\\
The Duke hath put on a religious life\\
And thrown into neglect the pompous court.
\end{verse_speech}

\verseline[Second Brother]{He hath.}

\begin{verse_speech}[Jaques]
To him will I. Out of these convertites\\
There is much matter to be heard and learned.\\
\stage{To \textsc{Duke}} You to your former honour I bequeath;\\
Your patience and your virtue well deserves it.\\
\stage{To \textsc{Orlando}} You to a love that your true faith doth merit.\\
\stage{To \textsc{Oliver}} You to your land, and love, and great allies.\\
\stage{To \textsc{Silvius}} You to a long and well-deserved bed.\\
\stage{To \textsc{Touchstone}} And you to wrangling, for thy loving voyage\\
Is but for two months victualled.—So to your pleasures.\\
I am for other than for dancing measures.
\end{verse_speech}

\verseline[Duke Senior]{Stay, Jaques, stay.}

\begin{verse_speech}[Jaques]
To see no pastime, I. What you would have\\
I'll stay to know at your abandoned cave.	
\end{verse_speech}

\exit{\textsc{Jaques}}

\begin{verse_speech}[Duke Senior]
Proceed, proceed. We'll begin these rites,\\
As we do trust they'll end, in true delights.
\end{verse_speech}

\stage{Dance}
\exit{all but \textsc{Rosalind}}

\cleardoubleoddpage
\thispagestyle{plain}
\vfill
\centerline{\LARGE Epilogue}
\vfill
\begin{prose_speech}[Rosalind]
  It is not the fashion to see the lady the epilogue, but it is no more unhandsome than to see the lord the prologue. If it be true that good wine needs no bush, 'tis true that a good play needs no epilogue. Yet to good wine they do use good bushes, and good plays prove the better by the help of good epilogues. What a case am I in then that am neither a good epilogue nor cannot insinuate with you in the behalf of a good play! I am not furnished like a beggar; therefore to beg will not become me. My way is to conjure you, and I'll begin with the women. I charge you, O women, for the love you bear to men, to like as much of this play as please you. And I charge you, O men, for the love you bear to women—as I perceive by your simpering, none of you hates them—that between you and the women the play may please. If I were a woman, I would kiss as many of you as had beards that pleased me, complexions that liked me, and breaths that I defied not. And I am sure as many as have good beards, or good faces, or sweet breaths will for my kind offer, when I make curtsy, bid me farewell.
  \end{prose_speech}
\exeunt{}
\vfill